Put \(c=1/e\).  We begin with two identities that hold under the one-sided support condition \(P\ll Q\).  If \(r=dP/dQ\), then \(E_Q[r]=1\).  Hence for any \(b\in\mathbb R\) and \(g_b(t)=f(t)-b(t-1)\),
\[
 D_{g_b}(P\|Q)=D_f(P\|Q)-b\{E_Q[r]-1\}=D_f(P\|Q). \tag{6}
\]
Also, using \(ec=1\),
\[
 B_{g_b}(e)=B_f(e)-b\{e(c-1)+(1-e)(0-1)\}=B_f(e). \tag{7}
\]

The exact cap class is always contained in the ball.  Indeed, if \(0\le r\le c\), convexity gives the chord inequality
\[
 f(r)\le \left(1-\frac r c\right)f(0)+\frac r c f(c).
\]
Integrating and using \(E_Q[r]=1\) yields
\[
 D_f(P\|Q)\le \left(1-\frac1c\right)f(0)+\frac1c f(c)
 =(1-e)f(0)+ef(1/e)=B_f(e). \tag{8}
\]
By \(\mathrm{thm:one\text{-}sided\text{-}bow\text{-}mixture\text{-}completeness}\), the cap condition characterizes \(\mathfrak L^{\to}_e(P)\), so
\[
 \mathfrak L^{\to}_e(P)\subseteq\mathfrak B^{\to}_{f,e}(P). \tag{9}
\]

Suppose \(f\in\mathfrak H_c\).  Choose \(b\) as in \(\mathrm{def:hinge\text{-}class}\) and write \(g=g_b\).  Then \(g=0\) on \([0,c]\), \(g>0\) on \((c,\infty)\), and thus \(g\ge0\).  Equations (6)--(7) give \(B_g(e)=0\) and
\[
 \begin{aligned}
 Q\in\mathfrak B^{\to}_{f,e}(P)
 &\Longleftrightarrow P\ll Q\text{ and }\int g(r)\,dQ=0\\
 &\Longleftrightarrow P\ll Q\text{ and }r\le c\quad Q\text{-almost surely}\\
 &\Longleftrightarrow Q\in\mathfrak L^{\to}_e(P).
 \end{aligned} \tag{10}
\]
The middle equivalence uses both pointwise nonnegativity and strict positivity above \(c\).  This proves universal equality.

For necessity, assume equality for every full-support binary \(P\).  We derive a supporting slope directly from convexity.  Define
\[
 b=\sup_{0\le s<1}\frac{f(1)-f(s)}{1-s}.
\]
This number is finite: the term with \(s=0\) is finite, and the secant-slope inequality for a convex function bounds every displayed left slope above by \(f(2)-f(1)\).  For \(t<1\), the definition of the supremum gives \(f(t)\ge f(1)+b(t-1)\).  For \(t>1\), convexity gives
\[
 \frac{f(t)-f(1)}{t-1}\ge\frac{f(1)-f(s)}{1-s}
 \quad\text{for every }s<1,
\]
and taking the supremum gives the same supporting inequality.  With \(g(t)=f(t)-b(t-1)\) and \(f(1)=0\), we therefore have
\[
 g(t)\ge0\quad(t\ge0),\qquad g(1)=0. \tag{11}
\]
Equations (6)--(7) show that replacing \(f\) by \(g\) changes neither the ball nor its radius.

For \(x\in(0,1)\), define
\[
 \lambda_x=\frac{1-x}{c-x},\qquad
 Q_x=\operatorname{Bernoulli}(\lambda_x),\qquad
 P_x=\operatorname{Bernoulli}(c\lambda_x).
\]
Both laws have full support.  On the success and failure atoms, respectively, the likelihood ratios \(dP_x/dQ_x\) are \(c\) and \(x\); the latter equality follows from
\[
 \frac{1-c\lambda_x}{1-\lambda_x}=x.
\]
Thus \(Q_x\in\mathfrak L^{\to}_e(P_x)\), and set equality implies ball feasibility.  In fact the chord bound is tight:
\[
 D_g(P_x\|Q_x)=B_g(e). \tag{12}
\]
To prove (12), suppose instead that the left side were strictly smaller.  For \(t>c\) put
\[
 \lambda(t)=\frac{1-x}{t-x},\qquad
 Q_t=\operatorname{Bernoulli}(\lambda(t)),\qquad
 P_t=\operatorname{Bernoulli}(t\lambda(t)).
\]
These remain full-support laws, with likelihood ratios \(t\) and \(x\).  The divergence
\[
 \lambda(t)g(t)+(1-\lambda(t))g(x)
\]
is continuous in \(t\) because \(g\) is continuous.  It would remain strictly below \(B_g(e)\) for all \(t>c\) sufficiently close to \(c\), putting \(Q_t\) in the ball although its first likelihood ratio exceeds the cap.  This contradicts binary set equality and proves (12).

Let \(x\uparrow1\).  Then \(\lambda_x\downarrow0\), and continuity together with \(g(1)=0\) shows that the left side of (12) tends to zero.  Therefore \(B_g(e)=0\).  Since
\[
 B_g(e)=\frac1c g(c)+\left(1-\frac1c\right)g(0),
\]
the nonnegativity in (11) and the positivity of both coefficients imply \(g(0)=g(c)=0\).  Convexity gives \(g(t)\le0\) between the endpoints, while (11) gives the reverse inequality; hence
\[
 g(t)=0\qquad(0\le t\le c). \tag{13}
\]
If \(g(t_0)=0\) for some \(t_0>c\), choose any \(x\in(0,1)\) and form the full-support binary pair whose two likelihood ratios are \(t_0\) and \(x\).  It violates the cap, but (13) and \(g(t_0)=0\) make its divergence zero, equal to \(B_g(e)\), contradicting set equality.  Thus \(g(t)>0\) for every \(t>c\), which is exactly \(f\in\mathfrak H_c\).

Universal equality trivially implies binary equality, and the preceding two directions prove the three-way equivalence.